\documentclass[10pt,twoside]{book}
\usepackage[utf8]{inputenc}
\usepackage[T1]{fontenc}
\usepackage{geometry}
\geometry{left=1.2cm,right=1.2cm,top=1.5cm,bottom=1.5cm}
\usepackage{amsmath,amssymb}
\usepackage{graphicx}
\usepackage{booktabs}
\usepackage{tabularx}
\usepackage{array}
\usepackage{xcolor}
\usepackage{titlesec}
\usepackage{setspace}
\usepackage{fancyhdr}
\usepackage{mdframed}
\usepackage{needspace}
\usepackage{enumitem}
\usepackage{hyperref}
\hypersetup{colorlinks=true,linkcolor=blue,urlcolor=blue}

\newcommand{\tag}[1]{\textsc{#1}}
\newcommand{\clk}[1]{\textit{[#1]}}
\newcommand{\stringitem}[1]{\textit{#1}}
\newcommand{\dusana}[1]{\begin{quote}\itshape #1 \\ --- Dusana of the Raven Road\end{quote}}

\titleformat{\chapter}{\Huge\bfseries}{\thechapter}{20pt}{}
\titleformat{\section}{\Large\bfseries}{\thesection}{15pt}{}
\setlength{\parindent}{0pt}
\setlength{\parskip}{0.8ex}

\newmdenv[linecolor=black,backgroundcolor=gray!5,innertopmargin=8pt,innerbottommargin=8pt,innerrightmargin=8pt,innerleftmargin=8pt]{infobox}

\begin{document}

\title{\Huge\textbf{The Wandering Ledger}\\ \large Tales of the Threshold, from the Eastern Roads}
\author{Told by the fires of a Viterran manor\\ \small Set down by Dusana of the Raven Road}
\date{Under the lamp, bread and salt --- and a debt that must be paid in stories}
\maketitle

\tableofcontents

\chapter*{Preface: The Borrowed Hearth}
\addcontentsline{toc}{chapter}{Preface: The Borrowed Hearth}

A ledger of stories is not a ledger of coin. It balances nothing --- but sometimes it buys time.

My people are the Tulkani. We build no temples. Our wagons carry the hearth. We do not own land; we \textit{borrow} it, season by season, until the road calls us elsewhere. The ground we rest on is never taken --- only kept warm until its true owners return. We call it \textit{borrowing the hearth}. The law calls it something else, but the law does not sleep under canvas.

This winter, the road did not call.

The passes closed early. Our wagons lost a wheel at the edge of Baronet Aldric's hunting forest --- a Viterran holding, marked by hedges and standing stones and a gate that had not been used in forty years. We borrowed a clearing. We lit a fire. We poured the first cup for the baronet, though he was not there.

He arrived at dawn, with three greyhounds and a constable.

"Trespass," he said. "Fine, imprisonment, or the gallows."

I offered him tea. He refused. I offered him bread and salt. He hesitated --- Viterran law respects hospitality --- and took a single bite.

"I cannot pay in coin," I said. "But I can pay in stories. One for each night we have borrowed your hearth. My family will repair the gate. I will speak."

He stared at me. Then he laughed --- a short, dry sound, like a branch breaking.

"You have been here seven nights. Seven stories. And if I am not satisfied, the eighth night, I hang the eldest."

I agreed. The eldest is my grandmother. She is ninety-two and has been ready to meet the Pale Shepherd for a decade. But I told the stories.

This book is the record of those seven nights. I have added three more --- for the nights the baronet came back, demanding more, pretending he was not enchanted.

The tales are not mine. They were told to me by rice farmers in Ayokha, by shrine maidens in Nihori, by silk merchants in Sihai, and by the small people --- the Lethai, the Aelinnel, the Aelaerem --- who live in the spaces between the great empires. I have changed names. I have hidden places. I have left out the ninth step, the ninth cup, the ninth word, because some taboos travel across all roads.

The baronet is still alive. My family is still borrowing his hearth. The gate is repaired.

Debt is a story we tell ourselves to make the world feel fair. The world is not fair. But the stories --- the stories are true enough.

\dusana{At the edge of Baronet Aldric's forest, under a lamp that flickers once for every untold tale. The ink is ash and mulberry sap.}

\chapter{The First Night: The Mirror in the Rice Paddy}

\section*{Told by}
Old Mei, a rice farmer in Ayokha who had never left her valley but knew the names of every spirit in the water.

\section*{The Tale}
In the monsoon season, the rice paddies become mirrors.

You see the sky above you, the clouds, the birds. But sometimes --- sometimes you see other things. A door that is not there. A face that does not belong to anyone you know. A version of yourself standing on the opposite bank, wearing different clothes.

A young farmer --- call him Heng --- saw his own reflection in the paddy at twilight. But the face was older, wearier, and it was smiling.

"You are going to sell the land," the reflection said.

"I am not," Heng said.

"You will. To pay the doctor. Your mother is sick."

Heng ran home. His mother was fine. She was cooking rice. He did not sell the land.

The next twilight, the reflection was younger. It was a child's face, his own, from before his father had died.

"You are going to sell the land," the child said.

"I am not," Heng said.

"You will. To pay the tax collector. The mandala demands its tithe."

Heng checked his ledgers. The tax collector was not due for three months. He did not sell.

The third twilight, the reflection was a stranger. It wore a silver mask with no eyeholes.

"You will sell the land," it said.

Heng was angry. "No."

The reflection reached out of the water and touched his foot. Its hand was cold. "Then your daughter will sell it. Or her daughter. The land does not belong to you. It belongs to the spirit who was here before the first seed was planted. You are only borrowing it. And the borrow is due."

Heng knelt. "What does the spirit want?"

"A story. A story that has never been told. Speak it into the water at the next full moon, and the debt is paid."

He did. He told the story of how his grandmother had found the land --- not by conquest, not by purchase, but by saving a child from drowning. The child was the spirit's own. The spirit had never known.

The water rippled. The reflection bowed. The land was his.

\section*{The Witch's Closing}
"Every field has a keeper. The keeper does not care about your deeds or your coin. It cares whether you know its name."

\section*{What Lurks in the Shadow of the Tale}
\begin{infobox}
\textbf{The Paddy Mirror} --- A spirit that lives in still water, visible only at twilight. It shows the future to those who do not ask, and the past to those who do not listen. Its touch is cold; its voice is the sound of water dripping.

\textbf{Tags / Mechanics:} \tag{SPIRIT} \tag{THRESHOLD} \tag{WISDOM}

\begin{itemize}
\item To speak with the spirit, a character must kneel at the water's edge at twilight and offer a story no one has heard. On success (\textit{Spirit + Lore}, DV 4), the spirit answers one question about the land.
\item On failure, the spirit shows a vision of the character's worst fear. The GM gains 1 SB, and the character marks 1 Fatigue.
\item \textbf{Clock: Borrowed Land [6]} -- ticks each time the land is used without acknowledgment. When full, the spirit demands a name, a memory, or a story. Refusal causes crops to fail for a season.
\end{itemize}

\textbf{Adventure Hook:} A village's rice harvest is failing. The farmers blame drought. The spirit says the land was sold without its permission. The party must find the missing deed --- or negotiate a new covenant.
\end{infobox}

\chapter{The Second Night: The Fox's Wedding}

\section*{Told by}
A Kitsune of no fixed name, met at a crossroads shrine in Nihori. She spoke in riddles and smelled of rain.

\section*{The Tale}
Foxes hold weddings in the rain.

You can see them if you are lost in a forest when the sun shines and the rain falls at the same time. A procession of lanterns, paper umbrellas, and small figures in silk. They do not invite humans. They do not forbid them --- but those who attend must pay a toll.

A woodcutter --- call him Kenji --- saw a fox's wedding in a clearing. He hid behind a cedar and watched.

The bride was a white fox with nine tails. The groom was a red fox with a silver bell around his neck. They exchanged vows in a language that sounded like wind through bamboo.

After the ceremony, the bride turned and looked directly at Kenji. "You have seen us," she said. "Now you must give us a gift."

Kenji had nothing. He gave them his axe.

The bride laughed. "An axe is for cutting. What do you wish to cut from your life?"

Kenji thought. "My debt."

The bride took the axe. "Then it is cut."

He returned home. His creditors had forgotten his name. His ledger was blank. He lived for forty years without owing a single coin.

When he died, the foxes came to his funeral. The bride laid the axe on his chest. "The debt is gone," she said. "But the axe remembers. Do not borrow again."

\section*{The Witch's Closing}
"A gift to a fox is a promise you cannot take back. Give only what you wish to lose forever."

\section*{What Lurks in the Shadow of the Tale}
\begin{infobox}
\textbf{The Fox's Wedding Procession} --- A spirit procession that appears during sun showers. It is neither malevolent nor kind. Those who witness it may offer a gift; the gift is always accepted, and the consequence is always proportional to the giver's hidden desire.

\textbf{Tags / Mechanics:} \tag{FAE} \tag{THRESHOLD} \tag{TRICK}

\begin{itemize}
\item The procession appears when a character is lost and the weather shifts abruptly. The character may choose to approach or hide.
\item Offering a gift (any possession of personal value) allows the character to "cut" one obligation: a debt, a geas, a relationship, or a memory. The GM removes that narrative thread from play.
\item The gift is consumed. The character cannot reclaim it.
\item \textbf{Clock: The Fox's Return [8]} -- ticks each time the character accepts a new debt or obligation. When full, the foxes return to collect "interest" --- a memory, a skill point, or a cherished relationship.
\end{itemize}

\textbf{Adventure Hook:} A noble has not aged in thirty years. She attended a fox's wedding and offered her name. Now the foxes haven"t returned to collect --- but she fears they will. The party must find a way to break the bargain without losing her identity.
\end{infobox}

\chapter{The Third Night: The Well That Did Not Sing}

\section*{Told by}
The daughter of a well-keeper in Ayokha, who learned the spirit's name from her mother and her mother's mother.

\section*{The Tale}
Every well has a singer. You cannot hear it unless you have drunk from the well and never wasted the water. The song is not music. It is the sound of the water remembering the sky.

One well --- in a village called Tham Luang --- stopped singing.

The villagers panicked. Without the song, the water grew bitter. The crops wilted. The children dreamed of drowning.

The well-keeper went down the rope. She found the spirit sitting at the bottom, its mouth sewn shut with silver thread.

"Who did this?" she asked.

The spirit could not answer. It pointed up.

She climbed out. She gathered the village elders. "Someone in this village has sewn the well's mouth. Confess, or the well will die."

No one confessed.

The well-keeper said, "Then I will cut the thread myself." She descended again, carrying a knife blessed by the god-king.

The spirit stopped her. It wrote in the mud: \textit{The thread is your own silence. You have not spoken the truth in twenty years.}

The well-keeper wept. She had been hiding a secret --- a death she had caused, a blame she had let fall on another. She climbed out. She confessed.

The thread dissolved. The well sang again.

The village did not forgive her. But they did not throw her out either. They built a new well, and she became its keeper. She sang to it every morning. The song told the truth. No one asked her what it said.

\section*{The Witch's Closing}
"The well sings your secrets. If you do not want to hear them, do not drink."

\section*{What Lurks in the Shadow of the Tale}
\begin{infobox}
\textbf{The Silent Well} --- A well whose spirit has been bound by an unspoken lie. It will not sing until the liar confesses. While silent, the water is bitter and causes nightmares.

\textbf{Tags / Mechanics:} \tag{LOCATION} \tag{TRUTH} \tag{CURSE}

\begin{itemize}
\item Drinking from the well forces a \textit{Wits + Resolve} test (DV 4) or the drinker must reveal a secret within the next day.
\item To restore the well, a character must speak a truth they have never told. The truth must be meaningful (GM discretion). The well then sings for one year.
\item \textbf{Clock: The Bitter Water [6]} -- ticks each time a character drinks from the well without confessing. When full, the drinker gains the Condition \textit{Silenced} (cannot speak in complete sentences until they confess or have the truth extracted by magic).
\end{itemize}

\textbf{Adventure Hook:} A village's wells have all gone silent. The keeper is dead, and no one knows the last secret she took to her grave. The party must find the truth --- or dig a new well.
\end{infobox}

\chapter{The Fourth Night: The Fhara Water-Clock}

\section*{Told by}
Suleiman the Oasis-Master, a Fhara caravaneer who carried a water-timer made of brass and glass. He never arrived late to any rendezvous.

\section*{The Tale}
The Fhara do not measure time by the sun. The sun lies. It moves differently in the desert than it does in the mountains, and it forgets to set when the wells are dry. Instead, they measure time by water.

A water-timer is a vessel with a small hole in the bottom. It fills at dawn, empties by noon, and is refilled three times before the evening prayer. There is a story that the first water-timer was not a timer at all, but a marriage contract.

A Fhara woman --- call her Fatima --- was promised to a man she did not love. Her father was a well-digger; her suitor was a spice merchant. The merchant offered seven camels and a chest of saffron. Fatima offered a counter-bargain.

"Let the water-timer decide," she said. "If it empties before I return from the deep well, I will marry you. If it still holds water when I climb out, you will leave this oasis and never return."

The merchant laughed. The deep well was thirty fathoms. She would be gone for hours. He agreed.

Fatima climbed down. She did not climb up. She sat at the bottom of the well, waiting. The water-timer dripped. One hour. Two. Three.

At the fourth hour, the water was still dripping. The merchant began to worry. At the fifth hour, he lowered a rope. Fatima tied it to a bucket of sand. He pulled up sand. She stayed below.

At the sixth hour, the water-timer was empty. The merchant called down, "The timer is dry! You have lost!"

Fatima laughed. "You filled the timer at dawn. Today is the summer solstice. The water evaporates faster than it drips. The timer has been telling lies since noon."

The merchant looked at the timer. The brass was hot. The water was gone. He had been tricked by his own contract.

He left. Fatima married a well-digger's son. They built a new water-timer that compensated for heat and cold. It is still used in that oasis, and every year on the solstice, they add a stone to a cairn --- one for each hour the timer is allowed to lie.

\section*{The Witch's Closing}
"Time is water. Water is a contract. Read the fine print before you let it drip."

\section*{What Lurks in the Shadow of the Tale}
\begin{infobox}
\textbf{The Water-Clock of the Solstice} --- A Fhara timepiece that measures water, not hours. It can be calibrated to tell different truths. On the solstice, it lies. Those who know its secrets can use the lie to break any timed contract.

\textbf{Tags / Mechanics:} \tag{ARTIFACT} \tag{TRICK} \tag{TRUTH}

\begin{itemize}
\item Once per year (GM's discretion), a character who understands the timer's mechanism may declare that a time-sensitive contract is invalid due to "evaporation error." The contract is voided; the other party suffers --2 dice to all actions for a day (confusion).
\item Using the timer for this purpose requires a \textit{Wits + Craft} test (DV 5). On a miss, the timer cracks and becomes useless until repaired (a downtime project).
\item \textbf{Clock: The Dripping Debt [4]} -- ticks each time the timer is used outside of the solstice. When full, the water spirit inside the timer awakens and demands a sacrifice (a memory, an oath, or a finger).
\end{itemize}

\textbf{Adventure Hook:} A merchant is trying to force a caravan master to pay a bond that comes due at noon. The caravan master has a water-timer --- but it was damaged in a storm. The party must repair it before the deadline, or find another way to prove the contract fraudulent.
\end{infobox}

\chapter{The Fifth Night: The Bridge of One Hundred Hands}

\section*{Told by}
An Aelinnel bell-tender who lived under a bridge and never learned to swim.

\section*{The Tale}
A bridge was built in the old style --- wooden planks, rope railings, and a single bronze bell at the midpoint. The bell would ring when the bridge was safe. It would not ring when the bridge was hungry.

A merchant --- call her Lin --- crossed the bridge at dusk. The bell did not ring. She waited. The bell did not ring. She walked anyway.

Halfway across, a hand reached up from the water. It was made of moss and bone. It grabbed her ankle.

"You did not wait for the bell," the hand said. "Now you owe a toll."

Lin had no coin. She offered her bracelet.

The hand refused. "Not coin. A hand. I have one hundred hands under this bridge. I need one more to reach the other side."

Lin looked down. In the water, she saw ninety-nine hands, reaching up, grasping. She saw the far bank, just a few steps away.

She knelt. "I will give you my hand. But first, let me write a letter to my daughter."

The hand released her. She wrote the letter. She tucked it into her sleeve. Then she held out her hand.

The bridge bell rang.

The hand withdrew. The other hands released. The water was still.

The merchant crossed. She never learned why the bell had changed its mind. But she left a new bell at the shrine on the far bank --- a small one, made of clay, with a single crack.

\section*{The Witch's Closing}
"A bridge without a bell is a mouth without teeth. It can still bite. Cross only when the bell rings --- or when you have nothing left to lose."

\section*{What Lurks in the Shadow of the Tale}
\begin{infobox}
\textbf{The Hand Bridge} --- A wooden bridge over a dark river. It is haunted by the spirits of those who drowned while the bridge was under construction. Each crossing requires a toll paid to the bell. If the bell does not ring, the toll is a hand --- or a story good enough to make the spirits reconsider.

\textbf{Tags / Mechanics:} \tag{LOCATION} \tag{THRESHOLD} \tag{DREAD}

\begin{itemize}
\item To cross safely, a character must ring the bell at the midpoint. If the bell rings, no toll is required.
\item If the bell is silent, the character must offer a hand (permanent injury: -1 die to all physical rolls) or tell a story of such power that the spirits weep. On a successful \textit{Spirit + Performance} test (DV 5), the spirits accept the story and the character crosses unharmed.
\item \textbf{Clock: The Hundredth Hand [8]} -- when full, the spirits have enough hands to pull the entire bridge into the water. Anyone on the bridge is drowned.
\end{itemize}

\textbf{Adventure Hook:} A bridge has collapsed, stranding two villages. The bell was stolen by a collector who did not believe the stories. The party must retrieve the bell --- or build a new bridge and appease the spirits with a hundred new stories.
\end{infobox}

\chapter{The Sixth Night: The Kuva's Road}

\section*{Told by}
An old Kuva woman at a caravanserai on the Way of Silk, who had walked from the Crimson Basin to the Inland Sea and back again. She smelled of cardamom and dust.

\section*{The Tale}
The Kuva do not build roads. They walk them. And when they walk, they leave behind a trail not of stones but of promises.

A young Kuva --- call him Dastan --- was tasked with delivering a message from the Satrapy of Dhahara to the court of the God-King in Ayokha. The message was sealed with brass and wax. He was not to open it.

He walked for forty days. On the forty-first, he met a Pereshi merchant who offered him a fortune to see the seal. Dastan refused.

"Then walk with me," the merchant said. "I will carry your pack."

Dastan agreed. They walked together for ten days. On the tenth, the merchant asked again. Dastan refused.

"Then let me carry your water," the merchant said.

Dastan agreed. They walked another ten days. On the tenth, the merchant asked again. Dastan refused.

"Then let me carry your debt," the merchant said. "I will pay whatever you owe."

Dastan stopped. "I owe nothing."

"Everyone owes something," the merchant said. "The road. The water. The ground you borrowed. You have been borrowing the world's patience for sixty days. Let me pay."

Dastan looked at the road ahead. He thought of his mother, who had died while he walked. He thought of the message he carried, which he would never read. He thought of the fortune he had refused.

"No," he said. "The debt is mine."

He walked on. The merchant did not follow.

When Dastan reached the God-King's court, he was forty days late. The message was a treaty --- an offer of peace between Dhahara and Ayokha. The offer had expired.

The God-King was angry. "You have brought us nothing but a dead promise."

Dastan knelt. "I have brought you forty days of walking. Forty days of refusing. Forty days of not borrowing what was not mine to take."

The God-King looked at him. "What is your name?"

"Dastan, son of no one."

"Then you will be the son of this court. Rise. The treaty is expired. But a promise kept is worth more than a treaty signed. We will make a new one."

Dastan became the God-King's messenger. He never carried a sealed letter again. He carried only his own word --- and it was enough.

\section*{The Witch's Closing}
"A borrowed road is not a stolen one. Walk it clean, and the road will remember. Walk it greedy, and the road will eat you."

\section*{What Lurks in the Shadow of the Tale}
\begin{infobox}
\textbf{The Unopened Seal} --- A message sealed with brass and wax. If carried without being opened for the length of a journey (at least 40 days), the carrier gains the ability to speak one truth that cannot be denied. The message itself becomes inert, but the carrier's word gains supernatural weight.

\textbf{Tags / Mechanics:} \tag{ARTIFACT} \tag{TRUTH} \tag{SACRIFICE}

\begin{itemize}
\item A character who carries a sealed message for 40 days without opening it may, upon delivery, declare one statement that listeners must accept as true (no roll). This can be used to end a dispute, force a confession, or reveal a lie.
\item The statement cannot be longer than a single sentence. After use, the character loses the ability to lie convincingly for a week (\textit{--2 dice to Deception}).
\item \textbf{Clock: The Road's Ledger [6]} -- ticks each time the character accepts a gift or favor while carrying the message. When full, the message opens itself, and the character forgets the message's contents. The GM gains 3 SB.
\end{itemize}

\textbf{Adventure Hook:} A courier has disappeared with a sealed treaty. The treaty is the only thing preventing a war. The party must find the courier --- but the courier has been walking for forty days, refusing to open the seal, and is now somewhere in the desert, guarded by a Pereshi merchant who wants the treaty for himself.
\end{infobox}

\chapter{The Seventh Night: The Bamboo Taboo}

\section*{Told by}
An Aelaerem hedge-witch who lived in a hollow tree and spoke only in odd numbers.

\section*{The Tale}
In a village at the edge of the Bamboo Forest, there was a rule: never sneeze seven times in a row.

If you sneezed once, you covered your mouth. Twice, you excused yourself. Thrice, you left the room. Four times, you drank ginger tea. Five times, you lay down. Six times, you called for the healer.

Seven times, you prayed.

A child --- call her Yuki --- did not know the rule. She sneezed seven times, and a shadow stood in the corner of her room.

"You have summoned me," the shadow said. "I am the Keeper of the Seventh Exhalation. I take what escapes."

Yuki did not understand. "What escapes?"

"Your name. Your shadow. Your last breath. You have none left."

The child ran to her grandmother. The grandmother knelt and said to the shadow, "Take mine instead."

The shadow considered. "Why?"

"Because I do not need it. I am old. I have sneezed seven times many years ago. I have been borrowing time ever since."

The shadow took the grandmother's name. The grandmother forgot who she was. But the child lived.

The village carved a new rule into the shrine gate: \textit{Count to six. Stop. The seventh sneeze is a door.}

\section*{The Witch's Closing}
"The body knows when to stop. The spirit knows when to ignore the body. Listen to both, or you will sneeze away something you cannot borrow back."

\section*{What Lurks in the Shadow of the Tale}
\begin{infobox}
\textbf{The Keeper of the Seventh Exhalation} --- A spirit that appears when a mortal sneezes seven times in a row. It claims the "seventh exhalation" --- which can be a name, a shadow, a memory, or a breath. It can be bargained with, but it always takes something.

\textbf{Tags / Mechanics:} \tag{SPIRIT} \tag{DREAD} \tag{TRUTH}

\begin{itemize}
\item A character who rolls a critical failure on an Endurance test may sneeze uncontrollably. The GM may declare that the character sneezes seven times, summoning the Keeper.
\item The Keeper demands a payment: a name (the character forgets a significant NPC), a shadow (the character loses their shadow for 1d4 days; --1 Position on Stealth), or a memory (lose 1 XP).
\item The character may offer a substitute: another character may volunteer a payment, or the party may offer a story (\textit{Spirit + Performance}, DV 5). On success, the Keeper takes the story and leaves.
\item \textbf{Clock: Borrowed Exhalations [4]} -- ticks each time the Keeper is summoned but not paid. When full, the Keeper claims the entire party's names for one night; they cannot introduce themselves, and no one remembers them until dawn.
\end{itemize}

\textbf{Adventure Hook:} A village has a plague of sneezing. The Keeper is collecting debts from the sick. The party must find a way to stop the sneezes --- or convince the Keeper to take a different toll.
\end{infobox}

\chapter*{Epilogue: The Eighth Night (Unrecorded)}
\addcontentsline{toc}{chapter}{Epilogue: The Eighth Night (Unrecorded)}

Baronet Aldric kept me for an eighth night. He said the seventh story was not enough. He said he needed to know what happened to the novice.

"She became the abbot," I said.

"And the stolen gold?"

"Was repaid."

"And the silences?"

"She learned to hear them. Then she went deaf, like the master. She said the silence was louder after forty years of bells."

The baronet stared into the fire. "Why did you stop at seven? Why not eight stories, or nine?"

"Because eight is the number of cups you pour for the living," I said. "Nine is for the Hollow. I do not owe the Hollow any stories. Not tonight."

He did not press. He is a Viterran. He knows the rule.

The next morning, he gave us a month's grace to repair the wagons. He did not ask for the ninth story. He did not ask for the stories I have added to this book.

But he did ask for my name. I gave him the one I use on the road. He nodded, wrote it in a ledger, and said, "If you pass this way again, I will have more firewood."

That is the closest a Viterran baronet comes to saying \textit{welcome}.

\dusana{Leaving the gate repaired, the debts unsettled, and the ninth bell un-rung. The wagon wheel holds --- for now.}

\chapter*{Postscript: A Letter Found in the Binding}
\addcontentsline{toc}{chapter}{Postscript: A Letter Found in the Binding}

To whoever reads this book after me:

I have told seven stories, and three more besides. But the baronet was right about one thing: a debt was left unpaid.

Not to him. To the road.

When I was young, I borrowed a tale from a Pereshi storyteller in the port of Zakov. He told it to me in exchange for a cup of tea. I promised to tell it to someone else --- to pass it along, as the Pereshi do. I never did. I forgot.

This winter, I remembered. The baronet's fire, the silence of the manor, the way his greyhounds listened with their ears forward --- it reminded me of that night in Zakov, with the salt wind and the storyteller's voice.

I have not told his tale here. It is not mine to write. But I have written a different tale --- the one about the Kuva and the sealed message --- and I have left a space for his.

The Pereshi storyteller's name was Kaveh. He had a scar on his left hand and a laugh that sounded like breaking glass. If you meet him, tell him the debt is acknowledged. He will know what it means.

And if you do not meet him --- tell the tale yourself. Borrow it. Keep it warm. Pass it on.

That is the only way to settle a debt to the Pereshi.

\dusana{In the wagon, as the snow begins to melt. The ink is running low. The road is calling.}

\chapter*{Appendix: Taboos and Customs of the Eastern Roads}
\addcontentsline{toc}{chapter}{Appendix: Taboos and Customs of the Eastern Roads}

The following customs appear in these tales. They are not exhaustive. They are not universal. They are simply true enough to be useful.

\begin{table}[h]
\centering
\begin{tabularx}{\textwidth}{p{4cm} p{4cm} X}
\toprule
\textbf{Custom} & \textbf{Culture} & \textbf{Consequence of Breaking} \\
\midrule
Never pour the ninth cup. & Ayokha, Aelaerem & The Hollow Bride may attend your feast. \\
Never sneeze seven times in a row. & Nihori (bamboo villages) & The Keeper claims your next exhalation. \\
Count to eight before opening a door that has been closed for a season. & Aelaerem (all) & The Hollow watches from the threshold. \\
Leave a gift for a fox's wedding. & Nihori (forest shrines) & The foxes remember. They always remember. \\
Do not drink from a silent well. & Ayokha (river valleys) & You must confess a secret or go thirsty. \\
Ring the bell before crossing a bridge. & Inland Sea (all) & The bridge demands a hand. \\
The seventh sneeze is a door. & Nihori, Aelaerem & \textit{Do not open it.} \\
Water-timer contracts expire at solstice. & Fhara (oasis courts) & The contract is void; the liar pays double. \\
A sealed message carried forty days cannot be opened. & Kuva (Way of Silk) & The word becomes truth; the carrier cannot lie. \\
A borrowed road is not a stolen one. & All & The road remembers; it may not let you pass again. \\
\bottomrule
\end{tabularx}
\end{table}

These are not laws. They are courtesies extended to the spirits who share the world. Violate them, and the spirits will not punish you --- they will simply stop being polite.

And impolite spirits are the most dangerous kind.

\chapter*{Epilogue of the Epilogue: What the Baronet Said}
\addcontentsline{toc}{chapter}{Epilogue of the Epilogue: What the Baronet Said}

The baronet read this book. He did not say whether he believed it.

But the next spring, a Tulkani wagon broke an axle on his land again. He did not send the constable. He sent a pot of tea.

"You owe me more stories," he said.

I told him I would write them when I returned from the south. I am going to Akilan, to the Sidhi coast, to the courts of the Pereshi. I have a debt to settle.

He nodded. "Then go. But leave a copy."

I did.

\dusana{At the crossroads, with a new wheel and an old promise. The ink is fresh. The road is south.}

\end{document}
